\section{Kernanforderungen der Aufgabenstellung}
Die Projektvorgabe definiert gewisse Grundvoraussetzung für die Umsetzung des Projekts.
Diese beginnen im Bereich des Datenmanagements mit der Partitionierung des Datensatzes in ein Trainings-, Validations- und Test-Set im Verhältnis von 60:20:20. 
In diesem Zuge erfolgt zudem die gezielte Handhabung fehlender Werte mittels Imputation (ersetzung fehlender Werte), sowie die Identifikation und Bereinigung von Ausreißern. 
Die Auswertung ist in zwei Teilbereiche gegliedert; eine statistische Analyse durch Tests und andere Statistische-Werkzeuge und eine deskriptive Aufbereitung durch Grafiken.
Die verwendeten Tests umfassen T-Tests und F-Tests.
Ebenso wurden vier lineare Modelle zur Vorhersage des Diamantenpreises spezifiziert, wo ein unterschiedlicher Ansatz bei der Merkmalsauswahl vorliegt.
Diese sind der Naiver-, Metrischer-, Vollmodell- und Domänenbasierter Ansatz. 
